\documentclass[9pt,aspectratio=169,xcolor=table]{beamer}

%~~~~~~~~~~~~~~~~~~~~~~~~~~~~~~~~~~~~~~~~~~~~~~~~~~~~~~~~~~~~~~~~~~~~~~~~~~~~~~
% Use roboto Font (recommended)
\usepackage[sfdefault]{roboto}
\usepackage[utf8]{inputenc}
\usepackage[T1]{fontenc}
%~~~~~~~~~~~~~~~~~~~~~~~~~~~~~~~~~~~~~~~~~~~~~~~~~~~~~~~~~~~~~~~~~~~~~~~~~~~~~~

%~~~~~~~~~~~~~~~~~~~~~~~~~~~~~~~~~~~~~~~~~~~~~~~~~~~~~~~~~~~~~~~~~~~~~~~~~~~~~~
% Define where theme files are located. ('/styles')
\usepackage{styles/fluxmacros}
\usefolder{styles}
% Use Flux theme v0.1 beta
% Available style: asphalt, blue, red, green, gray
% "red" is the closest built-in style to maroon; its exact colours are
% overridden with a true maroon palette just below, after the theme loads.
\usetheme[style=red]{flux}
%~~~~~~~~~~~~~~~~~~~~~~~~~~~~~~~~~~~~~~~~~~~~~~~~~~~~~~~~~~~~~~~~~~~~~~~~~~~~~~

%~~~~~~~~~~~~~~~~~~~~~~~~~~~~~~~~~~~~~~~~~~~~~~~~~~~~~~~~~~~~~~~~~~~~~~~~~~~~~~
% SKEE1033 maroon palette override.
% The flux theme's colortheme defines primaryLight/primary/secondary/tertiary
% as plain \definecolor calls (not \providecolor), so simply redefining them
% here, AFTER \usetheme, wins -- xcolor allows redefinition and every
% \setbeamercolor in beamercolorthemeflux.sty was already evaluated against
% these names, so the redefinition propagates everywhere automatically.
\definecolor{primaryLight}{HTML}{9B2242} % lighter maroon -- headers, accents
\definecolor{primary}{HTML}{7B1E3A}      % core maroon -- titles, block headers
\definecolor{secondary}{HTML}{3A5A6B}    % muted teal -- complementary accent
\definecolor{tertiary}{HTML}{8C6A3F}     % warm bronze -- example/third accent
%~~~~~~~~~~~~~~~~~~~~~~~~~~~~~~~~~~~~~~~~~~~~~~~~~~~~~~~~~~~~~~~~~~~~~~~~~~~~~~

%~~~~~~~~~~~~~~~~~~~~~~~~~~~~~~~~~~~~~~~~~~~~~~~~~~~~~~~~~~~~~~~~~~~~~~~~~~~~~~
% Extra packages
\usepackage{booktabs}
\usepackage{colortbl}
\usepackage{ragged2e}
\usepackage{amssymb}
\usepackage{multirow}
\usepackage{tikz}
\usetikzlibrary{arrows.meta, positioning, fit, backgrounds, calc, shapes.geometric}
\setbeamertemplate{itemize items}[square]
\usepackage[scaled=.92]{beramono}
\usepackage{listings}
\usepackage{array}
% Table striping: odd rows white, even rows very light gray.
\definecolor{stripe}{gray}{0.94}
\newcommand{\striped}{\rowcolors{1}{white}{stripe}}
%~~~~~~~~~~~~~~~~~~~~~~~~~~~~~~~~~~~~~~~~~~~~~~~~~~~~~~~~~~~~~~~~~~~~~~~~~~~~~~

%~~~~~~~~~~~~~~~~~~~~~~~~~~~~~~~~~~~~~~~~~~~~~~~~~~~~~~~~~~~~~~~~~~~~~~~~~~~~~~
% Code listing style (lightweight analogue of the book's skpython style).
% This course is Python-only, so Python is set as the lstlisting default via
% \lstset below -- every \begin{lstlisting} needs no language=... of its own.
\definecolor{codebg}{HTML}{FBF2F4}
\definecolor{codegreen}{rgb}{0,0.45,0}
\definecolor{codestring}{rgb}{0.58,0.1,0.2}
\lstdefinestyle{skpython}{
    basicstyle=\ttfamily\footnotesize,
    keywordstyle=\color{primary}\bfseries,
    commentstyle=\itshape\color{codegreen},
    stringstyle=\ttfamily\color{codestring},
    backgroundcolor=\color{codebg},
    breaklines=true,
    keepspaces=true,
    showspaces=false,
    showstringspaces=false,
    frame=single,
    rulecolor=\color{primary!40},
    numbers=none,
}
\lstset{language=Python, style=skpython}
%~~~~~~~~~~~~~~~~~~~~~~~~~~~~~~~~~~~~~~~~~~~~~~~~~~~~~~~~~~~~~~~~~~~~~~~~~~~~~~

%~~~~~~~~~~~~~~~~~~~~~~~~~~~~~~~~~~~~~~~~~~~~~~~~~~~~~~~~~~~~~~~~~~~~~~~~~~~~~~
% TikZ block styles for the computation-cycle and control-structure diagrams
\definecolor{cycfill}{HTML}{F3DCE4}
\definecolor{cycline}{HTML}{7B1E3A}
\definecolor{cycdofill}{HTML}{E8EEF0}
\definecolor{cycdoline}{HTML}{3A5A6B}
\tikzset{
  cycstep/.style = {font=\sffamily\small, align=center, rounded corners=2pt,
                     draw=cycline, fill=cycfill, line width=0.7pt,
                     minimum height=7mm, minimum width=20mm, inner sep=3pt},
  cycdo/.style   = {cycstep, fill=cycdofill, draw=cycdoline},
  flow/.style    = {-{Stealth[length=2mm]}, thick, draw=cycline},
}
%~~~~~~~~~~~~~~~~~~~~~~~~~~~~~~~~~~~~~~~~~~~~~~~~~~~~~~~~~~~~~~~~~~~~~~~~~~~~~~

%%%%%%%%%%%%%%%%%%%%%%%%%%%%%%%%%%%
%% SECTION SEPARATOR (ToC recap) %%
%%%%%%%%%%%%%%%%%%%%%%%%%%%%%%%%%%%
\AtBeginSection[]
{
{
    \setbeamercolor{background canvas}{bg=primaryLight!6}
    \begin{frame}[plain]
        \frametitle{Table of Contents}
        \tableofcontents[currentsection]
    \end{frame}
}
}
%%%%%%%%%%%%%%%%%%%%%%%%%%%%%%%%%%%

%~~~~~~~~~~~~~~~~~~~~~~~~~~~~~~~~~~~~~~~~~~~~~~~~~~~~~~~~~~~~~~~~~~~~~~~~~~~~~~
% Information
\title{Matrices and Circuit Equations}
\subtitle{\emph{Engineering Computation with Python} --- Chapter 9}
\author{Mun\'{i}m Zabidi}
\institute{\color{primary}Faculty of Electrical Engineering, Universiti Teknologi Malaysia}
\date{\today}
\titlegraphic{assets/logoutmwhite.png}
%~~~~~~~~~~~~~~~~~~~~~~~~~~~~~~~~~~~~~~~~~~~~~~~~~~~~~~~~~~~~~~~~~~~~~~~~~~~~~~

\begin{document}

\titlepage

\begin{frame}{Table of Contents}
    \tableofcontents
\end{frame}

%===============================================================================
\section{Why Circuits Need Matrices}
%===============================================================================

\begin{frame}{A Two-Loop Circuit}
    \leavevmode
    A circuit with resistors $R_1, R_2, R_3$ and sources $V_1, V_2$.
    Kirchhoff's voltage law around each loop gives two equations in the
    two unknown loop currents $I_1, I_2$:
    \[
    \begin{aligned}
    (R_1 + R_2)\,I_1 - R_2\,I_2 &= V_1, \\
    -R_2\,I_1 + (R_2 + R_3)\,I_2 &= V_2 .
    \end{aligned}
    \]

    \vspace{2mm}
    \begin{alertblock}{The problem}
        Solvable by substitution for two loops. Tedious for three.
        \alert{Impractical} for ten.
    \end{alertblock}
\end{frame}

\begin{frame}{The Same System, Written as a Matrix}
    \leavevmode
    All the information sits in a grid of coefficients and a column of
    source voltages:
    \[
    \underbrace{\begin{bmatrix} R_1 + R_2 & -R_2 \\ -R_2 & R_2 + R_3 \end{bmatrix}}_{\mathbf{A}}
    \underbrace{\begin{bmatrix} I_1 \\ I_2 \end{bmatrix}}_{\mathbf{x}}
    =
    \underbrace{\begin{bmatrix} V_1 \\ V_2 \end{bmatrix}}_{\mathbf{b}} .
    \]

    \vspace{2mm}
    \begin{block}{$\mathbf{A}\mathbf{x} = \mathbf{b}$}
        \textbf{A} holds what the circuit \alert{is}. \textbf{b} holds
        what \alert{drives} it. \textbf{x} holds what you want to
        \alert{find}.
    \end{block}

    \vspace{2mm}
    {\small This pattern --- one matrix equation standing in for many
    simultaneous equations --- is one of the most important in all of
    engineering.}
\end{frame}

\begin{frame}[fragile]{Solving It in Python}
    \leavevmode
    No need to invert $\mathbf{A}$ by hand, no need to rearrange the
    equations:

\begin{lstlisting}
R1, R2, R3 = 100.0, 220.0, 330.0    # ohms
V1, V2 = 12.0, 5.0                  # volts

A = np.array([[R1 + R2, -R2],
              [-R2,     R2 + R3]])
b = np.array([V1, V2])

I = np.linalg.solve(A, b)
print(f"I1 = {I[0]*1000:.3f} mA")
print(f"I2 = {I[1]*1000:.3f} mA")
\end{lstlisting}

    \texttt{I1 = 60.345 mA} \quad \texttt{I2 = 33.229 mA}
\end{frame}

\begin{frame}[fragile]{Check, Then Interpret}
    \leavevmode
    \begin{columns}[t]
        \begin{column}{.5\textwidth}
            \begin{block}{Check}
                Substitute back --- \alert{always}:
\begin{lstlisting}
print(A @ I)
# Output: [12.  5.]
\end{lstlisting}
                Matches $V_1, V_2$. This is the single most valuable
                habit in numerical linear algebra.
            \end{block}
        \end{column}
        \begin{column}{.5\textwidth}
            \begin{block}{Interpret}
                \small Both currents positive: both loop currents flow
                in the assumed direction. A negative result would just
                mean the true current flows the other way.

                \vspace{1mm}
                Tens of mA is reasonable for this circuit.
            \end{block}
        \end{column}
    \end{columns}

    \vspace{3mm}
    \begin{alertblock}{Predict before you run}
        If $R_2$ --- shared between the loops --- were made very large,
        what happens to $I_1$ and $I_2$? Predict, then set
        \texttt{R2 = 1e6} and check.
    \end{alertblock}
\end{frame}

%===============================================================================
\section{Matrix Operations in Detail}
%===============================================================================

\begin{frame}{What Does ``Solving'' Mean?}
    \leavevmode
    \[
    \begin{aligned}
    3x_1 + 2x_2 &= 12, \\
    x_1 + 4x_2 &= 10,
    \end{aligned}
    \qquad\Longleftrightarrow\qquad
    \begin{bmatrix} 3 & 2 \\ 1 & 4 \end{bmatrix}
    \begin{bmatrix} x_1 \\ x_2 \end{bmatrix}
    =
    \begin{bmatrix} 12 \\ 10 \end{bmatrix}.
    \]

    \vspace{2mm}
    Solving means finding the single vector $\mathbf{x}$ that satisfies
    \alert{all} the equations at once: $x_1 = 2.8$, $x_2 = 1.8$.

    \vspace{2mm}
    \begin{block}{What changes, what doesn't}
        The matrix form does not change the arithmetic; it changes the
        \alert{bookkeeping} --- a procedure written once works for two
        unknowns or two hundred.
    \end{block}
\end{frame}

\begin{frame}{Why Not Just Invert the Matrix?}
    \leavevmode
    Algebraically, $\mathbf{x} = \mathbf{A}^{-1}\mathbf{b}$ is correct.
    It is nevertheless the \alert{wrong} way to compute the answer:

    \vspace{2mm}
    \begin{columns}[t]
        \begin{column}{.5\textwidth}
            \begin{block}{It is slower}
                \small Forming the full inverse computes far more than
                the problem requires. Solving directly does roughly a
                third of the work for a large system.
            \end{block}
        \end{column}
        \begin{column}{.5\textwidth}
            \begin{block}{It is less accurate}
                \small Inversion introduces rounding at every step,
                which then propagates through the multiplication by
                $\mathbf{b}$.
            \end{block}
        \end{column}
    \end{columns}

    \vspace{3mm}
    \begin{alertblock}{The rule throughout this course}
        Use \texttt{np.linalg.solve(A, b)}, never
        \texttt{np.linalg.inv(A) @ b}.
    \end{alertblock}
\end{frame}

\begin{frame}[fragile]{When There Is No Solution}
    \leavevmode
    If one equation is a multiple of another:
    \[
    \begin{bmatrix} 1 & 2 \\ 2 & 4 \end{bmatrix}
    \]
    the two rows carry the same information --- no unique $\mathbf{x}$
    exists. This matrix is \alert{singular}; its determinant is zero.

\begin{lstlisting}
S = np.array([[1.0, 2.0],
              [2.0, 4.0]])
b = np.array([1.0, 2.0])

print(np.linalg.det(S))      # 0.0  -- singular
x = np.linalg.solve(S, b)    # raises LinAlgError
\end{lstlisting}

    \begin{block}{In a circuit problem}
        A singular \textbf{A} almost always means the equations were
        set up wrongly --- not that the circuit is impossible. The
        error is a useful early warning.
    \end{block}
\end{frame}

\begin{frame}{Matrix Multiplication, Reviewed}
    \leavevmode
    \[
    m = \begin{pmatrix} 1 & 2 \\ 3 & 4 \end{pmatrix}
    \]
    Computed element by element:

    \begin{center}
    \small
    \begin{tabular}{lll}
        \toprule
        \textbf{Result} & \textbf{Calculation} & \textbf{Value} \\
        \midrule
        $R_{0,0}$ & $(1{\times}1)+(2{\times}3)$ & $7$ \\
        $R_{0,1}$ & $(1{\times}2)+(2{\times}4)$ & $10$ \\
        $R_{1,0}$ & $(3{\times}1)+(4{\times}3)$ & $15$ \\
        $R_{1,1}$ & $(3{\times}2)+(4{\times}4)$ & $22$ \\
        \bottomrule
    \end{tabular}
    \end{center}

    \[
    m \times m = \begin{pmatrix} 7 & 10 \\ 15 & 22 \end{pmatrix}
    \]
\end{frame}

\begin{frame}{Transpose, Right Division, Left Division}
    \leavevmode
    \small
    \begin{columns}[t]
        \begin{column}{.33\textwidth}
            \begin{block}{Transpose}
                \scriptsize Flips a matrix over its main diagonal: rows
                become columns.
                \[
                a = \begin{pmatrix} 1 & 2 \\ 3 & 4 \end{pmatrix}
                \;\to\;
                a^T = \begin{pmatrix} 1 & 3 \\ 2 & 4 \end{pmatrix}
                \]
            \end{block}
        \end{column}
        \begin{column}{.33\textwidth}
            \begin{block}{Right division}
                \scriptsize Find $X$ solving $XA = B$:
                \[ X = BA^{-1} \]
            \end{block}
        \end{column}
        \begin{column}{.33\textwidth}
            \begin{block}{Left division}
                \scriptsize Find $X$ solving $AX = B$:
                \[ X = A^{-1}B \]
            \end{block}
        \end{column}
    \end{columns}

    \vspace{3mm}
    \begin{alertblock}{Order matters}
        Matrix multiplication is \alert{not} commutative:
        $BA^{-1} \neq A^{-1}B$. Right division and left division solve
        different equations.
    \end{alertblock}

    \vspace{2mm}
    {\scriptsize In practice, neither Python nor other languages compute
    $A^{-1}$ explicitly --- a numerically stable decomposition solves
    for $X$ directly.}
\end{frame}

\begin{frame}{Matrix Power}
    \leavevmode
    Raising a square matrix $A$ to an integer power $n$ means repeated
    \alert{matrix} multiplication:
    \[
    A^n = \underbrace{A \times A \times \cdots \times A}_{n \text{ times}}.
    \]

    \vspace{3mm}
    \begin{alertblock}{Not the same as element-wise power}
        $A^n$ (matrix power) and raising every \alert{element} of $A$
        to the power $n$ generally give different results, and are
        written differently in Python.
    \end{alertblock}
\end{frame}

%===============================================================================
\section{Matrix vs. Array Arithmetic}
%===============================================================================

\begin{frame}{Two Kinds of Arithmetic in NumPy}
    \leavevmode
    \begin{columns}[t]
        \begin{column}{.5\textwidth}
            \begin{block}{Matrix operations}
                \small Use \texttt{@}, or functions such as
                \texttt{np.dot()}, \texttt{np.matmul()},
                \texttt{np.linalg.solve()}.
            \end{block}
        \end{column}
        \begin{column}{.5\textwidth}
            \begin{block}{Element-wise (array) operations}
                \small Standard operators \texttt{+ - * /} act on each
                corresponding element, and work on arrays of any shape.
            \end{block}
        \end{column}
    \end{columns}

    \vspace{3mm}
    \begin{alertblock}{They agree on two operations}
        Addition (\texttt{+}) and subtraction (\texttt{-}) behave
        \alert{identically} whether you mean ``matrix'' or
        ``element-wise.'' Only multiplication, division, and power
        differ.
    \end{alertblock}
\end{frame}

\begin{frame}{Precedence of Matrix and Array Operators}
    \leavevmode
    As with ordinary arithmetic, NumPy follows an order of precedence:

    \begin{center}
    \small
    \begin{tabular}{cl}
        \toprule
        \textbf{Order} & \textbf{Operation} \\
        \midrule
        1 & Parentheses \\
        2 & Power \\
        3 & Transpose \\
        4 & Multiplication \& Division, left to right \\
        5 & Addition \& Subtraction \\
        6 & Colon \texttt{:} (slice) \\
        \bottomrule
    \end{tabular}
    \end{center}
\end{frame}

\begin{frame}{Algebra, Matrix Code, Array Code}
    \leavevmode
    \small
    \begin{center}
    \begin{tabular}{lll}
        \toprule
        \textbf{Algebra} & \textbf{Matrix (Python)} & \textbf{Array (Python)} \\
        \midrule
        $a + b$ & \texttt{a + b} & \texttt{a + b} \\
        $a - b$ & \texttt{a - b} & \texttt{a - b} \\
        $a \times b$ & \texttt{a @ b} & \texttt{a * b} \\
        $a / b$ (right div.) & \texttt{solve(a.T,b.T).T} & \texttt{a / b} \\
        $a \backslash b$ (left div.) & \texttt{solve(a, b)} & \texttt{b / a} \\
        $a^n$ & \texttt{matrix\_power(a,n)} & \texttt{a ** n} \\
        $a^T$ & \texttt{a.T} & \texttt{a.T} \\
        \bottomrule
    \end{tabular}
    \end{center}

    \vspace{3mm}
    \begin{alertblock}{A note on \texttt{np.matrix}}
        Legacy and not used in this course --- always \texttt{np.array}
        with \texttt{@}. Appendix B records the differences.
    \end{alertblock}
\end{frame}

%===============================================================================
\section{Worked Matrix Examples}
%===============================================================================

\begin{frame}[fragile]{Addition, Subtraction, Multiplication}
    \leavevmode
\begin{lstlisting}
A = np.array([[1, 2], [3, 4]])
B = np.array([[5, 6], [7, 8]])

A + B    # [[ 6  8] [10 12]]
A - B    # [[-4 -4] [-4 -4]]
A @ B    # [[19 22] [43 50]]
np.matmul(A, B)   # same as A @ B
\end{lstlisting}

    \begin{block}{Notice}
        Addition and subtraction look exactly like array arithmetic.
        Multiplication does not --- \texttt{@} is required for the
        true matrix product.
    \end{block}
\end{frame}

\begin{frame}[fragile]{Division and Power}
    \leavevmode
\begin{lstlisting}
# Right division: solving XA = B
np.linalg.solve(A.T, B.T).T
# [[-1.  2.] [-2.  3.]]

# Left division: solving AX = B
np.linalg.solve(A, B)
# [[-3. -4.] [ 4.  5.]]

# Matrix power
np.linalg.matrix_power(A, 2)
# [[ 7 10] [15 22]]

# Transpose
A.T   # [[1 3] [2 4]]
\end{lstlisting}
\end{frame}

\begin{frame}{Precedence in a Complex Expression}
    \leavevmode
    \small
    The rules from the previous section govern how Python evaluates a
    compound matrix expression:

    \begin{itemize}\itemsep2pt
        \item \textbf{Transpose before multiplication:} \texttt{A.T @ B}
        \item \textbf{Parentheses force order:} \texttt{(A + B) @ B}
        \item \textbf{Power before multiplication:}
              \texttt{matrix\_power(A,2) @ B}
        \item \textbf{Multiplication before addition:} \texttt{A @ B + B}
    \end{itemize}

    \vspace{2mm}
    \begin{block}{Worked check}
        \texttt{A.T @ B} $=
        \begin{pmatrix} 26 & 30 \\ 38 & 44 \end{pmatrix}$ \qquad
        \texttt{(A+B) @ B} $=
        \begin{pmatrix} 86 & 100 \\ 134 & 156 \end{pmatrix}$
    \end{block}
\end{frame}

\begin{frame}[fragile]{A Full Combination}
    \leavevmode
\begin{lstlisting}
# Power, transpose, addition, and division together
result = (np.linalg.matrix_power(A, 2) + A.T) \
          @ np.linalg.solve(B, A)
print(result)
\end{lstlisting}

    \texttt{[[-12.\ \ -7.] [-19. -10.]]}

    \vspace{3mm}
    \begin{block}{Reading it left to right}
        Power first (\texttt{matrix\_power(A,2)}), then transpose
        (\texttt{A.T}), then their sum, then \texttt{solve(B, A)}
        (left division), then the final matrix product.
    \end{block}
\end{frame}

%===============================================================================
\section{Complex Number Arithmetic}
%===============================================================================

\begin{frame}[fragile]{Pure Transpose vs.\ Conjugate Transpose}
    \leavevmode
    With complex-valued arrays, two different ``transposes'' exist:

\begin{lstlisting}
a = np.array([[1+1j, 2+2j], [3+3j, 4+4j]])

a.T          # pure transpose -- swaps rows/cols only
a.conj().T   # conjugate (Hermitian) transpose
\end{lstlisting}

    \begin{columns}[t]
        \begin{column}{.5\textwidth}
            \texttt{a.T}:
            \[
            \begin{pmatrix} 1{+}1j & 3{+}3j \\ 2{+}2j & 4{+}4j \end{pmatrix}
            \]
        \end{column}
        \begin{column}{.5\textwidth}
            \texttt{a.conj().T}:
            \[
            \begin{pmatrix} 1{-}1j & 3{-}3j \\ 2{-}2j & 4{-}4j \end{pmatrix}
            \]
        \end{column}
    \end{columns}

    \vspace{2mm}
    \begin{alertblock}{The difference}
        \texttt{.T} swaps rows and columns, leaving every real and
        imaginary part unchanged. \texttt{.conj().T} \alert{additionally}
        negates every imaginary part.
    \end{alertblock}
\end{frame}

%===============================================================================
\section{Summary and Practice}
%===============================================================================

\begin{frame}{What to Remember}
    \leavevmode
    \small
    \begin{itemize}\itemsep3pt
        \item Any set of simultaneous linear equations is a single
              matrix equation $\mathbf{A}\mathbf{x} = \mathbf{b}$.
        \item \texttt{np.linalg.solve(A, b)} solves it directly ---
              faster and more reliable than inverting $\mathbf{A}$.
        \item A \alert{singular} matrix (zero determinant) has no
              unique solution; \texttt{solve} raises \texttt{LinAlgError}.
        \item \texttt{@} is \alert{matrix} multiplication; \texttt{*} is
              \alert{element-wise}. Confusing them is one of the most
              common errors in numerical engineering code.
        \item \texttt{A.T} transposes; always check a solution by
              substitution (\texttt{A @ x} should reproduce \texttt{b}).
    \end{itemize}

    \vspace{3mm}
    \begin{alertblock}{The real skill}
        Setting the problem up correctly --- not solving it --- is
        where an engineer's judgement matters most.
    \end{alertblock}
\end{frame}

\begin{frame}[fragile]{Arithmetic Operator Drill}
    \leavevmode
    \small
    Before checking with Python, work these out by hand. Assume
    \texttt{a} is a scalar, \texttt{b} and \texttt{c} are arrays or
    matrices as appropriate:

    \begin{columns}[t]
        \begin{column}{.5\textwidth}
\begin{lstlisting}
a + 2 - 3
a + 2 * 3
2 * a**3 + a
a + b
np.dot(b, c)
c.T
np.dot(b, c.T)
b.T + a / 2 + 2
\end{lstlisting}
        \end{column}
        \begin{column}{.5\textwidth}
\begin{lstlisting}
np.dot(c.T, c)
c**(2**a)
np.dot(c, c)
c * c
b**b
a + b * b**c.T
(a + 2j).T
np.array([a+2j, a+3j]).T
\end{lstlisting}
        \end{column}
    \end{columns}
\end{frame}

\begin{frame}{Where This Is Going}
    \leavevmode
    \small
    This chapter answered how engineers write and solve systems of
    simultaneous equations compactly.

    \vspace{3mm}
    \begin{itemize}
        \item Chapter 10 returns to the two-loop circuit on a larger
              scale: reading many resistor/voltage configurations from
              a spreadsheet and writing the computed currents back.
        \item The same $\mathbf{A}\mathbf{x} = \mathbf{b}$ pattern
              reappears throughout engineering --- structural analysis,
              circuit networks, control systems.
    \end{itemize}

    \vspace{4mm}
    \begin{center}
        \large\color{primary}\textbf{Next: Measurements, Files and Spreadsheets}
    \end{center}
\end{frame}

\end{document}
